\section{Изводи и заключение}
\label{sec:conclusion}

Работата показа, че един октет е достатъчен, за да се реши дали нова TCP връзка е TLS
или некриптиран HTTP. Основата е в стандартите: първият запис на TLS клиента е
ръкостискане с ContentType 22, а HTTP методите се записват с главни латински букви.
Тестът, който минава през всичките 256 стойности, потвърждава, че двете множества не се
пресичат.

Самостоятелната реализация \code{octet} показва, че механизмът може да се изгради без
загуба на данни и без да се отслабва TLS. Прочетените октети се връщат в TLS
библиотеката чрез буфери в паметта. Ръкостискането се провежда от OpenSSL с минимална
версия TLS 1.2, без предоговаряне, без компресия и със строго договаряне на приложния
протокол. Един общ краен срок покрива и разпознаването, и ръкостискането, така че
прозорецът преди ръкостискането не дава на бавния клиент повече време от обикновения
TLS сървър. Тестовете минават на Windows и на Linux. Интеграционният тест откри реална
грешка в границата на заглавната част на HTTP, която тестът на ниво функция не можеше да
открие.

Цената на разпознаването е под една микросекунда на връзка в паметта и под
разделителната способност на измерването от край до край. При тези условия
разпознаването е практически безплатно и аргументът срещу него не може да е
производителността.

Аргументът срещу него, който остава, е сигурността. Сървър, който обслужва защитен и
незащитен транспорт на един порт, е съчетание, за което препоръките на IETF изрично
предупреждават заради атаките с понижаване на защитата. Разпознаването по първия октет
не създава този проблем, но и не го решава. Правилната употреба е некриптираният път
само да пренасочва към TLS, а отговорите по TLS да носят HSTS. В сегашния си вид
\code{octet} не прави това.

Естествените следващи стъпки са три. Първо, да се добавят пренасочването и HSTS, за да
стане сървърът годен за реална употреба. Второ, да се подложи на автоматизирано
подаване на случайни входни данни, най-вече анализаторът на HTTP и пътят на
разпознаването. Трето, да се измери на физическа машина с Linux и през истинска мрежа,
с клиент, който е по-бърз от сървъра, за да се стесни разделителната способност на
измерването от край до край.
